%%%%%%%%%%%%%%%%%%%%%%%%%%%%%%%%%%%%%%%%%
% eBook 
% LaTeX Template
% Version 1.0 (29/12/14)
%
% This template has been downloaded from:
% http://www.LaTeXTemplates.com
%
% Original author:
% Luis Cobo (luiscobogutierrez@gmail.com) with extensive modifications by:
% Vel (vel@latextemplates.com)
%
% License:
% CC BY-NC-SA 3.0 (http://creativecommons.org/licenses/by-nc-sa/3.0/)
%
%%%%%%%%%%%%%%%%%%%%%%%%%%%%%%%%%%%%%%%%%

%----------------------------------------------------------------------------------------
%	DOCUMENT CONFIGURATIONS AND INFORMATION
%----------------------------------------------------------------------------------------

\documentclass[oneside,11pt]{memoir} % Font size

\input{structure.tex} % Include the file that specifies the document structure and layout

\title{कौटिल्य की अर्थशास्त्र} % Book title
\author{कौटिल्य} % Author
\newcommand{\edition}{नयी टाइपिंग} % Book edition

%----------------------------------------------------------------------------------------

\begin{document}

%----------------------------------------------------------------------------------------
%	TITLE PAGE
%----------------------------------------------------------------------------------------

\thispagestyle{empty} % Suppress page numbering
\ThisCenterWallPaper{1}{kautilya.jpg} % Add the background image, the first argument is the scaling - adjust this as necessary so the image fits the entire page

ॐ | क्षत्रिय धर्म २१वीं सदी में
% \begin{tikzpicture}[remember picture,overlay]
% \node [rectangle, rounded corners, fill=white, opacity=0.75, anchor=south west, minimum width=3cm, minimum height=6cm] (box) at (-0.5,-10) (box){}; % White rectangle - "minimum width/height" adjust the width and height of the box; "(-0.5,-10)" adjusts the position on the page
% \node[anchor=west, color01, xshift=-1.5cm, yshift=-0.4cm, text width=2.9cm, font=\sffamily\scriptsize] at (box.north){\edition}; % "Text width" adjusts the wrapping width, "xshift/yshift" adjust the position relative to the white rectangle
% \node[anchor=west, color01, xshift=-1.5cm, yshift=-2cm, text width=2.9cm, font=\sffamily\bfseries\scshape\Large] at (box.north){\thetitle}; % "Text width" adjusts the wrapping width, "xshift/yshift" adjust the position relative to the white rectangle
% \node[anchor=west, color01, xshift=-1.5cm, yshift=-5cm, text width=2.9cm, font=\sffamily\bfseries] at (box.north){\theauthor}; % "Text width" adjusts the wrapping width, "xshift/yshift" adjust the position relative to the white rectangle
% \end{tikzpicture}

\newpage % Make sure the following content is on a new page

%----------------------------------------------------------------------------------------
%	TABLE OF CONTENTS
%----------------------------------------------------------------------------------------

\tableofcontents % Prints the table of contents

%----------------------------------------------------------------------------------------
%	INTRODUCTION SECTION
%----------------------------------------------------------------------------------------

\chapter*{परिचय} % Introduction chapter suppressed from the table of contents

\begin{quote}

\noindent\rule{7.5cm}{0.4pt}

ॐ गं गणपतये नमः 

\noindent\rule{7.5cm}{0.4pt}

'शारदा शारदाभौम्वदना। वदनाम्बुजे।

सर्वदा सर्वदास्माकमं सन्निधिमं सन्निधिमं क्रिया तू।'

श्रीं ह्रीं सरस्वत्यै स्वाहा।

ॐ ह्रीं ऐं ह्रीं सरस्वत्यै नमः।

\noindent\rule{7.5cm}{0.4pt}


बुद्धीन्द्रियमनःप्राणान्  जनानामसृजत्प्रभुः ।

मात्रार्थम् च भवार्थम् च आत्मनेकल्पनायच् ।।श्रीमद् भागवत् १०.८७.२।। 

अर्थात:

"श्रीभगवान् ने महाप्रलय के बाद सर्ग (स्रिष्टि) के प्रारम्भ मे आकाश, वायु, तेज, जल, पृथ्वी से मण्डित जगत् कि रचना की। जो जीव प्रलय के क्षण तक कैवल्य मोक्ष को नहीं प्राप्त हुए थे, यानि मुक्ति के बिना ही भगवान् में विलीन हो गए थे, ३१ नील १० ख़राब ४० अरब वर्षों तक, जो की मृत्यु लोक की गणना के हिसाब से है, भगवान् में विलीन रहे। उनके शुभ-अशुभ कर्मों को फल की ओर उन्मुख करने की भावना से और मोक्ष प्राप्ति के अनुकूल उनके जीवन में बल और वेग का संचार हो सके इसलिए भगवान् ने बुद्धि, इन्द्रिय, प्राण और अंतःकरण से पूर्ण जीवन प्रदान किया। " - श्री १४५ शंकराचार्य निश्चलानन्द सरस्वती, पूर्वाम्नाय गोवर्धन मठ , पूरी। 

\noindent\rule{7.5cm}{0.4pt}

सर्वधर्मान्परित्यज्य मामेकं शरणं व्रज।

अहं त्वा सर्वपापेभ्यो मोक्षयिष्यामि मा शुचः ।।श्रीमद्भगवद्गीता १८.६६।। 

अर्थात:

%यहाँ धर्म का अर्थ गुण जानना चाहिए। अतः सब गुण (जिसमे अच्छे बुरे सभी गुण आ जाते हैं) का त्याग कर मुझ एक की शरण ले। मैं तुझेसब पापों से मुक्त कर मोक्ष प्रदान करूँगा , चिंता न कर। 

\noindent\rule{7.5cm}{0.4pt}

\end{quote}

\section{यह पुस्तक क्यों?}

एक भाव जो भिन्न रास्तों से चल कर मुझ तक सर्वथा संचारित होता रहा है वह है लिखने का भाव। ये सब बातें लिखी थी हमारे पूर्वजों ने , उन पूर्वजों ने जो बात कहने का सामर्थ्य धारण करते थे। 
१. शब्द ब्रम्ह है। 
२. शब्द अक्षर से बने हैं जो अनादि हैं। 
३. अक्षर शिव डमरू की ध्वनि हैं। 
४. सर्वप्राचीन पुस्तक ऋग वेद है। 
५. वेद अपौरुषेय हैं। 
६. भागवत पुराण में वेदों की उत्त्पत्ति के सन्दर्भ में बात है १०.८७.२ श्लोक के आस पास। 

इस भाव से प्रेरित हो कर मै इस निश्चय पर पहुंचा हूँ की पुस्तक लिखने का सामर्थ्य एक सर्वोच्च सामर्थ्य है जो ब्रम्ह को उद्भाषित करने का एक मात्र सहज माध्यम है। 

I want to be ready for death before I die. Because, if I am not, then another birth is inevitable. I am a Hindu. I feel being blindfolded by a world that has exemplified arrogance in its approach towards nature. I am desperately touching every new object in my blurred sight with a hope that I will get a clue of whether I am on right path.

In this book I intend to articulate DOs and DON'Ts for myself with a hope that it could come handy for Indian youth. I envision this book to be a key to collective will of our glorious सनातन पुराण.
 
\footnote{There exists a desire in my heart to know what my ancestors knew. There exists a desire to have a guru to nudge me to optimal path for attaining freedom. Freedom from those to-do activities that lead to more to-dos [Influence by - NischalanandS 2019]. Before I begin delving deeper with complex questions, I would need to understand known answers}.

%----------------------------------------------------------------------------------------
%	BOOK PART
%----------------------------------------------------------------------------------------

\part{कौटिल्य अर्थशास्त्र for youth of Bharat in 21\textsuperscript{st} century}

%----------------------------------------------------------------------------------------
%	CHAPTER ZERO 
%----------------------------------------------------------------------------------------


\chapter{विनयाधिकारिक प्रथम अधिकरण}


पृथिव्या लाभे पालने च यवाणत्यर्थशास्त्राणि पूर्वाचार्यैः प्रस्तावितानि प्रायशस्तानि संहृत्यैक मिदमारथशास्त्रम् कृतम ।। 
तस्याम प्रकरणाधि करण समुद्द्येशः ।। २ ।। 

पृथ्वी के प्राप्त करने और प्राप्त की रक्षा करने के लिए जितने अर्थ-शास्त्र प्राचीन आचार्यों ने लिखे, प्रायः उन सबको ही समग्रीहीत, करके यह एक अर्थशास्त्र बनाया गया है। ।। १ ।।  सबसे प्रथम यह उसके प्रकरण और अधिकरणों का निरूपण किया जाता है। ।। २ ।। 

विद्यासमुद्देशः ।। ३ ।। वृद्ध संयोगः । । ४ । । इन्द्रियजयः । । ५ । ।  अमात्योत्पत्तिः । । ६ । ।  मंत्रिपुरोहितोत्पत्तिः । । ७ । । उपधाभिः शौचाशौच आत्यानम् । । 8 । ।  गूध्पुरुशोत्पत्तिः । । 9 । ।  गूध्पुरुष्प्रनिधिः । । 10। । स्वविषये कृत्यक्रुत्यपक्श्रक्षणम् । । 11। । परविष्ये कृत्यकृत्यपक्षो पग्रः । । 12। । मन्त्राधिकारः । । 13। । दूतप्रनिधिः । । 14। ।  राज्पुत्ररक्षणम् । । 15। । अव्रुद्ध्वृतम् । ।16। । अवरुद्धे च वृत्तिह्  । । 17। ।  राज्प्रनिधिः । । 18। । निशान्त्प्रनिधिः । । 19। । आत्म्रक्षित्कम् । । 20। ।  इति विनयाधिकारिकं प्रथम अधिकरणम् । । 21 । । 


%----------------------------------------------------------------------------------------
%	CHAPTER ONE
%----------------------------------------------------------------------------------------


\chapter{Sootra is seed | sow selectively}

My serious interest in thoughts of my ancestors began with संस्कृत. I was a teenager when my grandfather left his body. I remember the rituals performed by brahmins of my village beside one of its lakes. The sounds eminating from the chants performed by the brahmins was mesmerizing and it stayed in me like a seed waiting for right season. ``Where do I come from?'' is the question I began asking there on.

I am a Rajpoot. That segment among Rajpoots class of society that rule over 36 other classes.

%----------------------------------------------------------------------------------------
%	CHAPTER TWO
%----------------------------------------------------------------------------------------

\chapter{Rituals to sustain | take responsibilities}


%----------------------------------------------------------------------------------------
%	CHAPTER THREE
%----------------------------------------------------------------------------------------

\chapter{Give away fruits | to who deserves}



%----------------------------------------------------------------------------------------

\end{document}